\centerline{\bf \Huge Комбигеомные процессы}

\begin{flushright}
        \scriptsize{Эпизод 14. Seele, трон для души }
\end{flushright}

\problem{1}
На плоскости отмечено несколько точек общего положения. Докажите, что найдётся замкнутая несамопересекающаяся ломаная, множество вершин которой совпадает с множеством отмеченных точек.

\problem{2}
На координатной плоскости расположены несколько равных квадратов со сторонами, параллельными осям координат. Известно, что на плоскости нет точки, принадлежащей хотя бы трем квадратам. Докажите, что все квадраты можно раскрасить в три цвета так, чтобы никакие два квадрата одного цвета не имели общих точек.

\problem{3}
С четырехугольником $A B C D$ проделывают следующую операцию: одну из данных вершин меняют на точку, симметричную этой вершине относительно середины диагонали (концом которой она не является), обозначив новую точку прежней буквой. Эту операцию последовательно применяют к вершинам $A, B, C, D, A, B, \ldots$. Докажите, что в некоторые два разных момента времени получатся два равных четырехугольника.

\problem{4}
На плоскости даны 2023 точки общего положения, одна из них синяя, остальные красные. Докажите, что количество треугольников с вершинами в красных точках, содержащих синюю, чётно.

\problem{5}
На плоскости $n>1$ точек общего положения. Мельницей будем называть процесс, который начинается с прямой $\ell$, проходящей через отмеченную точку $P$, которую будем называть опорной. На каждом шаге прямая вращается против часовой стрелки вокруг опорной точки, пока не пройдёт через другую отмеченную точку $Q$. После этого $Q$ становится опорной точкой и процесс продолжается. Докажите, что можно выбрать отмеченную точку $P$ и прямую $\ell$, проходящую через $P$, так, чтобы мельница, начинающаяся из соответствующего положения, вращалась вокруг каждой точки бесконечное число раз.

\problem{6}
На плоскости отмечено $n$ точек общего положения. Будем говорить, что треугольник $A B C$ является \textit{красивым} для стороны $A B$, если площадь треугольника $A B C$ меньше площади треугольника $A B X$ для любой отмеченной точки $X$. Будем называть треугольник прекрасным, если он красивый хотя бы для двух своих сторон. Докажите, что найдётся не менее $\frac{1}{2}(n-1)$ прекрасных треугольников.